\section{Towards a Circular Economy?}
\label{sec:introduction}

For a long time now, leading natural scientists have warned that economic activities are pushing the Earth's ecosystems beyond their capacity to support the web of life \parencite{kendall1992,rockstrom2009,ripple2017,richardson2024}. Since a large part of the human pressure on the environment stems from the extraction and processing of materials and the emissions and deposition of waste products resulting from their use, environmentalists and green think tanks such as \textcite{ellenmacarthur2012} have advocated a transition to a ``circular economy'' where materials are reused or recycled back into production to the largest extent possible. Within the economics community, the ideas underlying the circular economy paradigm were anticipated long ago by writers such as \textcite{boulding1966} who criticized the existing ``cowboy economy'' with hardly any recycling and called for a new ``economics of spaceship Earth'' where ``Man must find his place in a cyclical ecological system which is capable of continuous reproduction of material form...'' \parencite[p.~7]{boulding1966}. Although mainstream economists have shown little interest in the circular economy concept, despite some attempts to demonstrate its relevance for environmental economics \parencite[e.g.,][]{sorensen2018,fullerton2022}, policymakers have been more receptive. In the European Union, a Circular Economy Act is due for adoption in late 2026, and in China circularity principles are being integrated into the country's Ecological Environment Code.\footnote{Information on these policy initiatives may be found at \url{https://environment.ec.europa.eu/strategy/circular-economy_en} and \url{https://enviliance.com/regions/east-asia/cn/report_15527}. \textcite{mhatre2021} offer a broader survey of circular economy initiatives in the European Union.}

Using an extended Ramsey growth model with a mining industry and a waste treatment sector and pollution externalities in production and consumption, this paper shows that at some stage during the process of economic development it becomes socially optimal to move from a so-called linear economy with no recycling of materials to a circular economy with an increasing degree of recycling. The transition to a circular economy is driven by increasing natural resource scarcity and growing pressures on the environment during the development process. However, since recycling requires the use of resources which may at some point generate more harm than good to the environment, as already pointed out by \textcite{baumol1977}, there is an optimal degree of recycling well below a hundred percent. The paper contributes to environmental and public macroeconomics by deriving a set of Pigouvian taxes and subsidies that will ensure a socially optimal transition to a circular economy with an optimal degree of recycling. Importantly, we find that Pigouvian taxes on polluting activities are not sufficient for optimality; they must be supplemented by Pigouvian subsidies to waste collection and recycling. We also simulate a numerical version of our model to illustrate how technical progress in various parts of the economy affect the transition to a circular economy and the optimal degree of recycling and how much welfare is likely to be gained from the optimal environmental policy relative to a scenario with a laissez-faire policy.

The literature on the circular economy is overwhelming. For example, the survey by \textcite{kirchherr2023} identifies no less than 221 different definitions of ``circular economy''. Academic fields such as industrial ecology and business management have seen a particularly strong interest in the concept, and this literature is too large to be reviewed here, but surveys may be found in \textcite{winans2017,merli2018,centobelli2020,ferasso2020}. Our paper adds to a relatively small environmental economics literature on recycling. An early contribution was made by \textcite{smith1972} who focused on the reuse of household waste. \textcite{schultze1973} illustrated how the recycling of raw materials could ameliorate the exhaustion of non-renewable resources, and \textcite{lusky1975,lusky1976} studied the allocation of household time between work in the labour market and recycling activity, showing in the 1975 paper how the optimal amount of recycling might be secured through a tax on consumption. \textcite{hoel1978} analyzed the optimal path of economic development and the role of recycling when natural resource extraction harms the environment. The simplicity of his model generated the stark conclusion that resource extraction and recycling will never take place simultaneously, whereas the present analysis finds that the two activities can go on at the same time. The more recent papers by \textcite{divita2001,divita2007} investigate how endogenous technical change driven by R\&D may affect the recycling of waste and thereby consumer welfare, and \textcite{pittel2010} set up a Ramsey-type model of exogenous growth with recycling of waste to study how the optimal level of recycling may be implemented through government subsidies. Like the present paper, \textcite{andersen2007} makes the point that the policy problems discussed within the circular economy paradigm can be tackled via the classical Pigouvian policy instruments emphasized in conventional environmental economics. The contributions mentioned above did not focus on explaining the transition from a linear to a circular economy. That transition was the focus of \textcite{sorensen2018}, but in his simple model a Pigouvian tax on natural resource extraction was sufficient to secure the optimal degree of recycling. In the much richer model developed here a broader set of public finance instruments is shown to be needed.

Our paper combines elements of neoclassical economics and ecological economics. We undertake a utilitarian welfare analysis balancing marginal social costs and benefits of economic activities as in neoclassical economics, but inspired by ecological economics we explicitly account for the materials balance principle and the first and second laws of thermodynamics, and our quantitative analysis allows for limited possibilities of substitution between natural resources and man-made capital when raw material use falls to low levels.

The paper is structured as follows. Section~\ref{sec:model-setup} describes the assumptions on tastes and technologies in our model. On this basis, Section~\ref{sec:first-best-allocation} characterizes the first-best optimal allocation of resources and analyzes the conditions that will trigger an optimal transition from a linear economy with no recycling of materials to a circular economy with waste treatment and an increasing degree of recycling. Section~\ref{sec:market-implementation} shows how the government can implement the optimal allocation in a market economy by an appropriate choice of Pigouvian taxes and subsidies. Section~\ref{sec:quantitative-analysis} sets up and calibrates a numerical version of the model and simulates the optimal evolution of the economy over time, illustrating how the transition to a circular economy is affected by the rates of technical progress in various sectors of the economy. We also use the numerical model to estimate the welfare loss incurred if the government fails in various ways to implement the optimal tax and subsidy policy. Section~\ref{sec:conclusion} discusses the insights from and the limitations of the analysis.
